\begin{table}[htbp]
\centering
\captionsetup{font=small}
\caption{Experimental Attributes, Levels, and Coding Used in the Discrete Choice Experiment}
\small
\begin{tabularx}{0.9\textwidth}{@{} l Y Y Y @{}}
\toprule
\textbf{Attribute} & \textbf{Levels} & \textbf{Coding} & \textbf{Description} \\
\midrule
Vaccine Delays &
0, 1, 3, 6 months (Walk-in to 6 mo) &
0, 1, 3, 6 &
Waiting time before the vaccine becomes available. \\

Vaccine Origin &
Domestic; Imported &
0 = Domestic, 1 = Imported &
Indicates whether the vaccine was produced locally or overseas. \\

Vaccine Efficacy &
50\%, 70\%, 95\% &
0.50, 0.70, 0.95 &
Probability of protection against infection. \\

Side Effects &
Mild; Moderate; Severe &
0, 1, 2 &
Ordered severity of expected vaccine side effects. Higher values indicate more severe reactions. \\

Monetary Incentives (RMB) &
50, 200, 800 &
50, 200, 800 &
Monetary compensation offered upon receiving the vaccine. \\
\bottomrule
\end{tabularx}

\vspace{0.25em}
\captionsetup{font=footnotesize}
\caption*{\textit{Notes:} Attribute codings reflect the numerical values entered into the econometric model. Side-effect severity is treated as an ordinal variable.}
\end{table}

